\documentclass[10pt]{mypackage}

\usepackage[uprightgreeks,light]{kpfonts}
\renewcommand{\coloneq}{\coloneqq}
\usepackage{homework}
%\usepackage{notes}

%\usepackage[ backend=bibtex, style = alphabetic, sorting=ynt ]{biblatex}
%\addbibresource{  }

\usepackage{parskip}

\fancyhf{}
\fancyhead[R]{Avinash Iyer}
\fancyhead[L]{Algebraic Topology: Homework 21 Resubmission}
\fancyfoot[C]{\thepage}

\setcounter{secnumdepth}{0}

\begin{document}
\RaggedRight
\begin{problem}[Problem 1]
  Let $\left( X,A \right)$ be a pair of spaces.
  \begin{enumerate}[(a)]
    \item Prove that $H_0\left( X,A \right) = 0$ if and only if $A$ intersects each path component of $X$ non-trivially.
    \item Prove that $H_1\left( X,A \right) = 0$ if and only if the map $H_1\left( A \right)\rightarrow H_1\left( X \right)$ is surjective and each path component of $X$ contains at most one path component of $A$.
  \end{enumerate}
\end{problem}
\begin{solution}\hfill
  \begin{enumerate}[(a)]
    \item Suppose $H_0\left( X,A \right) = 0$. From the long exact sequence, we then have that the map $H_0\left( A \right)\xrightarrow{i_{\ast}} H_0\left( X \right)$ is surjective, meaning that the homology classes $\left[ \alpha \right]\in H_0\left( A \right)$ corresponding to path components of $A$ generate $H_0\left( X \right)$ upon inclusion into $X$, so the collection of path components of $A$ surjects onto the path components of $X$, meaning the intersection of $A$ with every path component of $X$ is nontrivial.

      In the reverse direction, if $A$ intersects every path component of $X$ nontrivially, then there are homology class representatives $\left[ \alpha \right]$ in $H_0\left( A \right)$ corresponding to each path component of $H_0\left( X \right)$ via the inclusion of $A$ into $X$, so $H_0\left( A \right)$ surjects onto $H_0\left( X \right)$, meaning $H_0\left( X,A \right) = 0$ by exactness.
    \item Suppose $H_1\left( X,A \right) = 0$. Then, in the long exact sequence, we observe that we have the segment
      \begin{align*}
        H_1\left( A \right) \xrightarrow{i_{\ast}} H_1\left( X \right)\xrightarrow{j_{\ast}}0,
      \end{align*}
      meaning that $i_{\ast}$ is surjective. Furthermore, we have
      \begin{align*}
        0 \xrightarrow{\partial} H_0\left( A \right)\xrightarrow{i_{\ast}}H_0\left( X \right),
      \end{align*}
      meaning that $i_{\ast}$ is injective. If we let $\left[ \alpha_i \right]\in H_0\left( A \right)$ denote the homology classes in $H_0\left( A \right)$ corresponding to the path components of $A$, then the inclusion of these homology classes into $H_0\left( X \right)$ leads to an injection of path components of $A$ into path components of $X$.

      In the reverse direction, if every path component of $X$ contains at most one path component of $A$, then the homology classes in $H_0\left( A \right)$ corresponding to the path components inject into the homology classes of $H_0\left( X \right)$ corresponding to the path components. Therefore, we have $H_0\left( A \right)\xrightarrow{i_{\ast}}H_0\left( X \right)$ is an injective map. If $H_1\left( A \right)\xrightarrow{i_{\ast}}H_1(X) $ is surjective, then by the long exact sequence, we have
      \begin{align*}
        H_1\left( A \right)\xrightarrow{i_{\ast}}H_1\left( X \right)\xrightarrow{j_{\ast}}H_1\left( X,A \right)\xrightarrow{\partial} H_0\left( A \right)\xrightarrow{i_{\ast}}H_0\left( X \right)
      \end{align*}
      with $i_{\ast}$ at the $H_1$ level surjective and $i_{\ast}$ at the $H_0$ level injective. Thus, by exactness, we must have $H_1\left( X,A \right) = 0$.
  \end{enumerate}
\end{solution}
\end{document}
